% Компилируется и pdfLaTeX, и XeLaTeX/LuaLaTeX (дважды)
\documentclass[aspectratio=169,9pt]{beamer}

\usetheme{Madrid}
\usecolortheme{seahorse}
\setbeamertemplate{navigation symbols}{}
\setbeamertemplate{footline}[frame number]
\setbeamertemplate{itemize items}[circle]
\setbeamertemplate{bibliography item}{\insertbiblabel}

\usepackage{iftex}
\ifPDFTeX
  % pdfLaTeX (например, Overleaf по умолчанию)
  \usepackage[T2A]{fontenc}
  \usepackage[utf8]{inputenc}
  \usepackage[russian]{babel}
\else
  % XeLaTeX / LuaLaTeX
  \usepackage{fontspec}
  \setmainfont{DejaVu Serif}
  \setsansfont{DejaVu Sans}
  \setmonofont{DejaVu Sans Mono}
\fi

\usepackage{amsmath,amssymb}
\usepackage{booktabs}
\usepackage{tikz}
\usetikzlibrary{arrows.meta,positioning,calc,decorations.pathreplacing}

\newcommand{\E}{\mathbb{E}}
\newcommand{\KL}{\mathrm{KL}}
\newcommand{\fx}{\mathbf{f}^{\mathbf{X}}}
\newcommand{\bx}{\mathbf{X}}

\title[Functional Variational BNNs]{Functional Variational Bayesian Neural Networks}
\subtitle{Sun, Zhang, Shi, Grosse --- ICLR 2019 (arXiv:1903.05779)}
\author{Chirkov George}
\institute{Bayesian multimodeling, МФТИ}
\date{2026}

\begin{document}

%% 1 ---------------------------------------------------------------
\begin{frame}
  \titlepage
  \note{Бюджет времени: около 10 минут, 15 слайдов. Титульный слайд 0:15.}
\end{frame}

%% 2 ---------------------------------------------------------------
\begin{frame}{Почему вывод в пространстве весов работает плохо}
  Bayesian neural networks (BNN): prior $p(\mathbf w)$, likelihood $p(\mathcal D\mid\mathbf w)$, posterior $p(\mathbf w\mid\mathcal D)$ невычислим $\Rightarrow$ вариационный вывод, например Bayes by Backprop (BBB): $q(\mathbf w)=\mathcal N(\boldsymbol\mu,\mathrm{diag}\,\boldsymbol\sigma^2)$.
  \vspace{0.4em}
  \begin{block}{Проблемы VI в пространстве весов}
    \begin{itemize}
      \item Prior-ы неосмысленны. Веса связаны с вычисляемой функцией сложным образом; на практике $\mathcal N(0,I)$ берут для удобства.
      \item Posterior трудно приблизить. Между весами сильные зависимости, а $q$ (почти) факторизована.
      \item Деградация при росте ёмкости. Чем больше сеть, тем сложнее вывод и больше description length у posterior. На $y=x^3$ (20 точек) предсказания BBB резко меняются с размером сети.
    \end{itemize}
  \end{block}
  Разрыв: GP-prior-ы кодируют богатую структуру (гладкость, периодичность), но масштабируются как $O(N^3)$; BNN масштабируются, но их prior-ы неинтерпретируемы.
  \note{Рис. 1 статьи: BBB против fBNN на x^3. Упомянуть устно. 1:00.}
\end{frame}

%% 3 ---------------------------------------------------------------
\begin{frame}{Идея: неопределённость над функциями, а не над весами}
  \begin{alertblock}{Идея}
    Вариационный вывод напрямую над распределениями по функциям: prior --- стохастический процесс, приближённый posterior --- стохастическая нейросеть.
  \end{alertblock}
  \vspace{0.3em}
  \begin{columns}[T]
    \begin{column}{0.5\textwidth}
      \textcolor{structure}{Что мы получаем}
      \begin{itemize}
        \item структурные prior-ы: GP (периодичность, гладкость) и \emph{implicit} процессы (кусочные функции)
        \item явные сэмплы функций из posterior (Thompson sampling, entropy search)
        \item предсказания стабильны при росте сети
        \item масштабируемость: данные входят только через ожидаемый log-likelihood (мини-батчи)
      \end{itemize}
    \end{column}
    \begin{column}{0.5\textwidth}
      \textcolor{structure}{Вклад статьи}
      \begin{enumerate}
        \item Теорема: KL между процессами = supremum маргинальных KL по конечным наборам точек
        \item Functional ELBO (fELBO) как игра типа GAN; sampling-based версия
        \item Оценка градиентов через spectral Stein gradient estimator (SSGE)
        \item Эксперименты: экстраполяция, UCI-регрессия, bandits, Bayesian optimization
      \end{enumerate}
    \end{column}
  \end{columns}
  \note{0:45.}
\end{frame}

%% 4 ---------------------------------------------------------------
\begin{frame}{Что именно мы хотим посчитать}
  Данные: $\mathcal D=\{(\mathbf x_i,y_i)\}_{i=1}^n$, likelihood $p(\mathcal D\mid f)=\prod_i p(y_i\mid f(\mathbf x_i))$.

  \medskip
  Prior: стохастический процесс $p$ над функциями $f:\mathcal X\to\mathcal Y$ (GP или implicit процесс без плотности).

  \medskip
  Вариационное семейство: стохастическая сеть с шумом $\xi$, общим для всех входов
  \[
    f(\mathbf x)=g_\phi(\mathbf x,\xi),\qquad \xi\sim p(\xi)
    \quad(\text{epistemic uncertainty}).
  \]
  Цель: максимизировать functional ELBO
  \[
    \mathcal L(q)=\E_q\big[\log p(\mathcal D\mid f)\big]-\KL[q\,\|\,p].
  \]
  \begin{block}{Сравнение с пространством весов}
    $\mathcal L_q=\E_q[\log p(\mathcal D\mid\mathbf w)]-\KL[q(\mathbf w)\|p(\mathbf w)]$.
  \end{block}
  Сложности: (i) $\KL[q\|p]$ между процессами не имеет вида $\int q\log\frac qp$; (ii) плотность $q$ над $f$ невычислима; (iii) prior может быть implicit.
  \note{1:00.}
\end{frame}

%% 5 ---------------------------------------------------------------
\begin{frame}{Как измерить расстояние между двумя процессами}
  Процесс характеризуется согласованными конечномерными маргиналами (теорема Колмогорова о продолжении: exchangeability и consistency).

  \begin{block}{Теорема 1 (Functional KL)}
    Для двух стохастических процессов $P$ и $Q$
    \[
      \KL[P\,\|\,Q]=\sup_{n\in\mathbb N,\ \bx\in\mathcal X^n}\KL[P_{\bx}\,\|\,Q_{\bx}].
    \]
  \end{block}
  Почему: $\sigma$-алгебра из теоремы о продолжении порождена цилиндрическими множествами, зависящими лишь от конечного числа точек.

  \medskip
  Следствие. Для measurement set $\bx$ и $\mathcal L_{\bx}(q):=\E_q\sum_{i}\log p(y_i\mid f(\mathbf x_i))-\KL[q(\fx)\|p(\fx)]$
  \[
    \mathcal L(q)=\inf_{n,\ \bx\in\mathcal X^n}\mathcal L_{\bx}(q).
  \]
  $\Rightarrow$ антагонистическая игра с нулевой суммой: один игрок выбирает сеть $q$, противник --- measurement set (как в GAN).
  \note{1:00.}
\end{frame}

%% 6 ---------------------------------------------------------------
\begin{frame}{В каких точках смотреть на функции}
  \begin{columns}[T]
    \begin{column}{0.48\textwidth}
      \textcolor{structure}{Adversarial} (как в GAN)
      \[
        \max_{q}\ \min_{\bx\in\mathcal X^m}\ \mathcal L_{\bx}(q)
      \]
      \begin{itemize}
        \item худшие точки лежат рядом с обучающими данными
        \item ничто не заставляет $q$ следовать структуре prior в других областях $\Rightarrow$ плохое обобщение
      \end{itemize}
    \end{column}
    \begin{column}{0.52\textwidth}
      \textcolor{structure}{Sampling-based} (используется в fBNN)
      \[
        \max_{q}\ \E_{\bx^M\sim c}\ \mathcal L_{\bx^M,\bx^{D_s}}(q)
      \]
      \begin{itemize}
        \item $\bx^{D_s}$: случайный мини-батч обучающих точек
        \item $\bx^M$: $M$ случайных точек из интересующей области ($c$)
        \item согласуем сеть с prior и рядом с данными, \emph{и} там, где нужны предсказания
      \end{itemize}
    \end{column}
  \end{columns}

  \vspace{0.6em}
  \centering
  \begin{tikzpicture}[x=0.85cm,y=0.8cm]
    \draw[->,thick] (-0.3,0) -- (10.6,0) node[right] {$\mathbf x$};
    \foreach \x in {2.5,2.8,3.2,3.6,6.4,6.8,7.2,7.5}
      \fill[red] (\x,0.25) circle (2.2pt);
    \foreach \x in {0.4,1.3,2.0,4.6,5.2,5.9,8.4,9.1,9.8}
      \draw[blue,thick] (\x-0.1,0.15) -- (\x+0.1,0.35) (\x-0.1,0.35) -- (\x+0.1,0.15);
    \draw[decorate,decoration={brace,amplitude=4pt,mirror}] (0,-0.2) -- (10,-0.2)
      node[midway,below=4pt] {\scriptsize область сэмплирования $[x_{\min}-\tfrac d2,\ x_{\max}+\tfrac d2]$, $d=x_{\max}-x_{\min}$};
    \node[red,anchor=south] at (3.0,0.75) {\scriptsize обучающие входы};
    \node[blue,anchor=south] at (8.6,0.75) {\scriptsize случайные measurement points};
  \end{tikzpicture}
  \note{0:50. Картинка схематична; она следует схеме сэмплирования из раздела экспериментов статьи.}
\end{frame}

%% 7 ---------------------------------------------------------------
\begin{frame}{Когда objective --- настоящая нижняя граница}
  \begin{block}{Теорема 2 (Lower bound)}
    Если measurement set $\bx$ содержит все обучающие входы $\bx^{D}$, то
    \[
      \mathcal L_{\bx}(q)=\log p(\mathcal D)-\KL\big[q(\fx)\,\|\,p(\fx\mid\mathcal D)\big]\ \le\ \log p(\mathcal D).
    \]
  \end{block}
  При $\bx=\{\bx^D,\bx^M\}$: максимизация $\mathcal L_{\bx}(q)$ $\Leftrightarrow$ минимизация KL до истинного posterior в этих точках (то же тождество ELBO, что на лекции).

  \begin{block}{Следствие 3 (Consistency)}
    Если истинный posterior $p(f\mid\mathcal D)$ --- GP, $\mathcal Q$ --- все GP, $M>1$ и $\mathrm{supp}(c)=\mathcal X$, то оптимумы и adversarial, и sampling-based задач равны $p(f\mid\mathcal D)$.
  \end{block}
  Оговорка: для конечного measurement set, не содержащего все обучающие входы, objective --- лишь \emph{приближение} (не гарантированная граница для $\log p(\mathcal D)$); функциональный KL может быть даже бесконечным.
  \note{0:50.}
\end{frame}

%% 8 ---------------------------------------------------------------
\begin{frame}{Градиент KL, когда у $q$ нет плотности}
  \small
  Likelihood-член вычислим, а плотность $q_\phi(\fx)$ --- нет. По репараметризации $f=g_\phi(\mathbf x,\xi)$:
  \[
    \nabla_\phi\KL[q_\phi(\fx)\|p(\fx)]
    =\underbrace{\E_q[\nabla_\phi\log q_\phi(\fx)]}_{=0}
    +\E_\xi\Big[\nabla_\phi\fx\,\big(\nabla_{\mathbf f}\log q(\fx)-\nabla_{\mathbf f}\log p(\fx)\big)\Big].
  \]
  Произведение якобиана на вектор считается backprop-ом; остаётся оценить score-функции $\nabla_{\mathbf f}\log q$ и $\nabla_{\mathbf f}\log p$.

  \begin{block}{Spectral Stein gradient estimator (SSGE)}
    Оценивает $\nabla\log q(\mathbf x)$ \emph{implicit} распределения только по сэмплам:
    \[
      \nabla_{x_i}\log q(\mathbf x)=-\sum_{j=1}^{\infty}\big[\E_q\nabla_{x_i}\psi_j(\mathbf x)\big]\psi_j(\mathbf x),
    \]
    $\psi_j$: собственные функции ядра $k$ (Mercer), приближаются методом Nystr\"om; сумма обрезается, $\E_q$ заменяется оценкой Monte Carlo.
  \end{block}
  SSGE применяется для $\nabla\log q$ всегда и для $\nabla\log p$, когда prior implicit; для GP-prior $\nabla\log p$ считается точно (с шумом $K+\gamma^2I$ для устойчивости).
  \note{1:00.}
\end{frame}

%% 9 ---------------------------------------------------------------
\begin{frame}{Собираем всё вместе: как обучается fBNN}
  \centering
  \begin{tikzpicture}[
      box/.style={draw,rounded corners,align=center,minimum height=0.8cm,inner sep=3pt,font=\footnotesize},
      arr/.style={-{Stealth},thick},
      node distance=0.9cm and 0.7cm]
    \node[box] (batch) {мини-батч\\$\bx^{D_s}$};
    \node[box,below=0.35cm of batch] (meas) {случайные точки\\$\bx^M\sim c$};
    \node[box,right=1.1cm of $(batch)!0.5!(meas)$,fill=blue!10] (net) {стохастическая сеть\\$f=g_\phi(\mathbf x,\xi)$};
    \node[box,right=0.8cm of net] (f) {$\mathbf f^{D_s},\ \mathbf f^{M}$\\(сэмплы функций)};
    \node[box,right=0.8cm of f,yshift=0.75cm] (ll) {ожидаемый\\log-likelihood};
    \node[box,right=0.8cm of f,yshift=-0.75cm] (kl) {KL-член\\градиенты SSGE};
    \node[box,right=0.9cm of $(ll)!0.5!(kl)$,fill=green!10] (upd) {обновление $\phi$\\через backprop};
    \draw[arr] (batch) -- (net); \draw[arr] (meas) -- (net);
    \draw[arr] (net) -- (f);
    \draw[arr] (f) -- (ll); \draw[arr] (f) -- (kl);
    \draw[arr] (ll) -- (upd); \draw[arr] (kl) -- (upd);
  \end{tikzpicture}

  \medskip
  \raggedright
  Objective, максимизируемый на каждой итерации:
  \[
    \frac1{|\mathcal D_s|}\sum_{(\mathbf x,y)\in\mathcal D_s}\E_{q_\phi}\log p(y\mid f(\mathbf x))\;-\;\lambda\,\KL\big[q(\mathbf f^{\mathcal D_s},\mathbf f^M)\,\|\,p(\mathbf f^{\mathcal D_s},\mathbf f^M)\big].
  \]
  \begin{itemize}
    \item В принципе $\lambda=1/|\mathcal D|$; так как KL по малому числу точек --- лишь нижняя оценка функционального KL, большее $\lambda$ помогает против переобучения. На практике $\lambda=1/|\mathcal D_s|$ (тогда по теореме 2 это честная нижняя граница $\log p(\mathcal D_s)$).
    \item $q$ во всех экспериментах: независимые гауссовы веса $\mathcal N(\boldsymbol\mu,\mathrm{diag}\,\boldsymbol\sigma^2)$; гиперпараметры GP сначала подбираются, затем фиксируются.
  \end{itemize}
  \note{1:00.}
\end{frame}

%% 10 --------------------------------------------------------------
\begin{frame}{Может ли сеть экстраполировать по структуре prior}
  \begin{columns}[T]
    \begin{column}{0.5\textwidth}
      \textcolor{structure}{Периодическая структура (GP-prior)}
      \begin{itemize}
        \item 20 точек из $[-2,-0.5]\cup[0.5,2]$, $y=2\sin(4x)+\varepsilon$
        \item BNN: 5 скрытых слоёв $\times$ 500 нейронов
        \item prior-ы: ядра RBF и PER$+$RBF
        \item measurement set: все обучающие точки $+$ 40 случайных из $[-5,5]$
      \end{itemize}
      Результат: fBNN $\approx$ точный GP с тем же ядром (чуть меньшая неопределённость). BBB не может даже подогнать обучающие данные.
    \end{column}
    \begin{column}{0.5\textwidth}
      \textcolor{structure}{Implicit prior-ы}
      \begin{itemize}
        \item кусочно-постоянные и кусочно-линейные функции
        \item 40 точек из $[0,0.2]\cup[0.8,1]$; сеть с \emph{tanh}, 2$\times$100
      \end{itemize}
      Результат: сэмплы из posterior имеют резкие, примерно кусочные структуры, несмотря на tanh, но не отражают всё разнообразие объяснений.

      \medskip
      Toy-пример $y=x^3$: предсказания fBNN согласованы при росте сети; предсказания BBB резко меняются.
    \end{column}
  \end{columns}
  \vspace{0.5em}
  \begin{block}{Вывод}
    Экстраполяция следует из \emph{functional prior}, а не из архитектуры (ReLU-сети не кодируют периодичность).
  \end{block}
  \note{1:00. При наличии времени упомянуть пример с CO2 на Mauna Loa из приложения.}
\end{frame}

%% 11 --------------------------------------------------------------
\begin{frame}{Сравнение с BBB и Noisy K-FAC на регрессии (test RMSE)}
  \small
  \textcolor{structure}{Малые датасеты UCI} (GP-prior с подобранными гиперпараметрами; split 90/10, 10 повторов)
  \begin{center}
  \begin{tabular}{lccc}
    \toprule
    Датасет & BBB & Noisy K-FAC & fBNN \\
    \midrule
    Boston   & 3.171 $\pm$ 0.149 & 2.742 $\pm$ 0.125 & \textbf{2.378} $\pm$ 0.104 \\
    Concrete & 5.678 $\pm$ 0.087 & 5.019 $\pm$ 0.127 & \textbf{4.935} $\pm$ 0.180 \\
    Energy   & 0.565 $\pm$ 0.018 & 0.485 $\pm$ 0.023 & \textbf{0.412} $\pm$ 0.017 \\
    Wine     & 0.643 $\pm$ 0.012 & \textbf{0.637} $\pm$ 0.011 & 0.673 $\pm$ 0.014 \\
    Yacht    & 1.174 $\pm$ 0.086 & 0.979 $\pm$ 0.077 & \textbf{0.607} $\pm$ 0.068 \\
    \bottomrule
  \end{tabular}
  \end{center}
  \textcolor{structure}{Большие датасеты} (1 скрытый слой, 100 нейронов; GP-prior с NKN-ядрами; 5 запусков)
  \begin{center}
  \begin{tabular}{lrcc}
    \toprule
    Датасет & $N$ & BBB & fBNN \\
    \midrule
    Naval        & 11\,934  & 1.6E-4 & \textbf{1.2E-4} \\
    Protein      & 45\,730  & 4.331  & \textbf{4.326} \\
    Video Memory & 68\,784  & 1.879  & \textbf{1.858} \\
    Video Time   & 68\,784  & 3.632  & \textbf{3.007} \\
    GPU          & 241\,600 & 21.886 & \textbf{19.50} \\
    \bottomrule
  \end{tabular}
  \end{center}
  \note{0:50. По log-likelihood картина смешанная на Video Memory/Time (BBB лучше); упомянуть, если спросят.}
\end{frame}

%% 12 --------------------------------------------------------------
\begin{frame}{Bandits и Bayesian optimization: зачем нужна неопределённость}
  \textcolor{structure}{Contextual bandits} (Riquelme et al., 2018), Thompson sampling, prior = multi-task GP. Относительный накопленный regret (Uniform $=100$), 6 датасетов, 10 запусков:
  \begin{center}
  \small
  \begin{tabular}{lcc}
    \toprule
    Метод & Mean rank $\downarrow$ & Mean value $\downarrow$ \\
    \midrule
    fBNN $1\times 50$   & 4.7  & 41.9 \\
    fBNN $1\times 500$  & \textbf{3.8}  & 41.3 \\
    fBNN $3\times 500$  & 4.2  & \textbf{40.9} \\
    MultitaskGP         & 4.3  & 41.7 \\
    Dropout             & 5.5  & 41.7 \\
    BBB $1\times 50$    & 10.8 & 52.7 \\
    BBB $1\times 500$   & 13.7 & 66.2 \\
    Uniform             & 16   & 100 \\
    \bottomrule
  \end{tabular}
  \end{center}
  \begin{itemize}
    \item fBNN стабилен при более глубоких и широких сетях; у BBB качество падает при переходе от 50 к 500 нейронам.
    \item Bayesian optimization (Max-value Entropy Search): fBNN даёт явные сэмплы функций; сопоставим с RBF random features на RBF-функциях и лучше на ArcCosine и Matern12.
  \end{itemize}
  \note{0:50.}
\end{frame}

%% 13 --------------------------------------------------------------
\begin{frame}{Что получилось и где слабые места}
  \begin{columns}[T]
    \begin{column}{0.5\textwidth}
      \textcolor{structure}{Выводы}
      \begin{itemize}
        \item $\KL$ между процессами = supremum конечных маргинальных KL
        \item fELBO, sampling-based measurement sets и SSGE дают практичную схему обучения
        \item prior-ы над функциями (GP, implicit) $\Rightarrow$ осмысленная экстраполяция, надёжная неопределённость, устойчивость к размеру сети
        \item конкурентно на UCI-регрессии, сильный результат на bandits
      \end{itemize}
    \end{column}
    \begin{column}{0.5\textwidth}
      \textcolor{structure}{Ограничения}
      \begin{itemize}
        \item конечные measurement sets: лишь приближение, $\lambda$ подбирается эвристически
        \item сэмплы posterior для implicit prior-ов недостаточно разнообразны
        \item гиперпараметры GP подбираются заранее и фиксируются
        \item $q$ по-прежнему факторизованная гауссиана по весам; масштабируемость SSGE в больших размерностях не исследована
      \end{itemize}
    \end{column}
  \end{columns}
  \note{0:30.}
\end{frame}

%% 14 --------------------------------------------------------------
\begin{frame}{Литература}
  \footnotesize
  \begin{thebibliography}{99}
    \bibitem{fbnn} Sun S., Zhang G., Shi J., Grosse R. \emph{Functional variational Bayesian neural networks.} ICLR 2019, arXiv:1903.05779.
    \bibitem{bbb} Blundell C., Cornebise J., Kavukcuoglu K., Wierstra D. \emph{Weight uncertainty in neural network.} ICML 2015.
    \bibitem{ssge} Shi J., Sun S., Zhu J. \emph{A spectral approach to gradient estimation for implicit distributions.} ICML 2018.
    \bibitem{neal} Neal R. M. \emph{Bayesian Learning for Neural Networks.} PhD thesis, University of Toronto, 1995.
    \bibitem{gpml} Rasmussen C. E., Williams C. K. I. \emph{Gaussian Processes for Machine Learning.} MIT Press, 2006.
    \bibitem{nkfac} Zhang G., Sun S., Duvenaud D., Grosse R. \emph{Noisy natural gradient as variational inference.} ICML 2018.
    \bibitem{bandits} Riquelme C., Tucker G., Snoek J. \emph{Deep Bayesian bandits showdown.} ICLR 2018.
    \bibitem{matthews} Matthews A. G. de G., Hensman J., Turner R., Ghahramani Z. \emph{On sparse variational methods and the KL divergence between stochastic processes.} AISTATS 2016.
    \bibitem{ncp} Hafner D. et al. \emph{Reliable uncertainty estimates in deep neural networks using noise contrastive priors.} arXiv:1807.09289, 2018.
    \bibitem{pbp} Hern\'andez-Lobato J. M., Adams R. \emph{Probabilistic backpropagation for scalable learning of Bayesian neural networks.} ICML 2015.
  \end{thebibliography}
\end{frame}

%% 15 --------------------------------------------------------------
\begin{frame}{Два вопроса к аудитории}
  \begin{enumerate}
    \setlength{\itemsep}{1.2em}
    \item Почему задание prior над функциями (а не над весами) делает предсказания fBNN стабильными при росте сети?
    \item Зачем в measurement set входят случайные точки из всей области определения, а не только обучающие входы?
  \end{enumerate}
  \vfill
  \centering\Large Спасибо за внимание!
  \note{%
    Ответ 1: задача функционального вывода (prior над f, KL между процессами) определена без привязки к архитектуре, поэтому увеличение сети лишь расширяет вариационное семейство и не меняет саму задачу.
    Ответ 2: KL-член --- единственная часть, которая заставляет сеть следовать структуре prior; на одних обучающих входах он ничего не говорит про области, где нужна экстраполяция, поэтому добавляют точки из области предсказаний.}
\end{frame}

\end{document}
